%=============================================================================!
% 3.7  FRICTION AND THE LOSS OF ENERGY                                        !
%=============================================================================!
\section{Friction and the loss of energy}
\label{sec:en-friction}

A box shoved across a floor slides to a stop, and a real bead on a real
wire never quite climbs back to the height it was released from. Both
lose energy to friction. This section measures the loss and finds the
rule that says how fast the total energy changes when part of the force
is not conservative.

\subsection*{Sliding friction}

When one surface slides over another, each exerts on the other a
friction force along the surfaces, and the force on each points against
the direction in which that surface slides over the other. For many pairs of dry surfaces, the size of this force depends only
weakly on the speed and on the area in contact, and is roughly
proportional to the normal force $F_{\mathrm{n}}$ with which the surfaces
are pressed together,
\begin{equation*}
   F_{\mathrm{fr}} = \mu_{\mathrm{k}}\,F_{\mathrm{n}} .
\end{equation*}
The number $\mu_{\mathrm{k}}$, the coefficient of sliding friction,
depends on the two surfaces and is found by measurement. This rule is an empirical description of limited
accuracy, and we use it only to see how energy is lost.

A box of \SI{5}{\kilogram} slides across a level floor at
\SI{2}{\metre\per\second}, with $\mu_{\mathrm{k}} = 0.3$. Vertically the
box does not accelerate, so the normal force balances the weight,
$F_{\mathrm{n}} = mg$, and the friction is $0.3\times5\times9.80665 =
\SI{14.71}{\newton}$, pointing against the motion. The weight and the
normal force are perpendicular to the motion and do no work, so over a
distance $d$ the work on the box is that of friction, $-\mu_{\mathrm{k}}
mg\,d$. The box stops when this has removed all of its kinetic energy,
\begin{align*}
   0 - \tfrac12 m v_0^2 &= -\mu_{\mathrm{k}} m g\,d ,\\
   d &= \frac{v_0^2}{2\mu_{\mathrm{k}}g}
      && \text{dividing both sides by $-\mu_{\mathrm{k}}mg$}\\
   &= \frac{2^2}{2\times0.3\times9.80665} = \SI{0.680}{\metre} .
\end{align*}
The mass cancels again: a heavier box has proportionally more kinetic
energy to lose, and it presses harder on the floor and feels a
proportionally larger friction.

The \SI{10}{\joule} of kinetic energy the box had is not destroyed. The
rubbing surfaces warm up, and Chapter~12 shows that warmth is the hidden
kinetic and potential energy of the jiggling atoms of the box and the
floor. Energy has passed from the motion of the box as a whole, which we
follow, into motions of its atoms, which we do not.

\subsection*{Down a rough ramp}

A block of \SI{2}{\kilogram} slides from rest down a straight ramp
\SI{3}{\metre} long, inclined at \ang{30}, with $\mu_{\mathrm{k}} = 0.2$
(\cref{fig:en-ramp}). Three forces act. The line perpendicular to the
ramp is the vertical turned through the same \ang{30} as the ramp is
turned from the level, so it makes an angle of \ang{30} with the weight,
$mg$ straight down. As for the sledge of \cref{sec:en-constant}, the
weight therefore has a component $mg\cos\ang{30}$ into the ramp and a
component $mg\sin\ang{30}$ at right angles to that, down the ramp. The
normal force balances the first of these, because the block does not
accelerate into or away from the ramp, so $F_{\mathrm{n}} = mg\cos\ang{30}
= \SI{16.986}{\newton}$. Friction acts up the ramp with size
$\mu_{\mathrm{k}}F_{\mathrm{n}} = 0.2\times16.986 = \SI{3.397}{\newton}$.

\begin{figure}[htb]
\centering
\begin{tikzpicture}[line cap=round, line join=round, scale=1.15,
   force/.style={-{Stealth[length=5pt]}, line width=1.0pt},
   part/.style={-{Stealth[length=4pt]}, line width=0.6pt, dashed}]
   % the ramp, rising to the left at 30 degrees
   \fill[black!10] (0,0) -- (5.2,0) -- (0,3.0) -- cycle;
   \draw[line width=0.6pt] (0,0) -- (5.2,0) -- (0,3.0);
   \draw[black!50, line width=0.4pt] (4.3,0) arc[start angle=180,
      end angle=150, radius=0.9];
   \node at (3.95,0.2) {\small \ang{30}};
   % the block, its centre on the line through the middle of its face
   \begin{scope}[shift={(2.25,1.70)}, rotate=-30]
      \draw[line width=0.6pt, fill=white] (-0.45,0) rectangle (0.45,0.6);
      \coordinate (C) at (0,0.3);
   \end{scope}
   % weight and its two parts
   \draw[force] (C) -- ++(0,-1.35) coordinate (W)
      node[below] {\small $mg$};
   \draw[part] (C) -- ++(-30:0.675) node[right] {\small $mg\sin\ang{30}$};
   \draw[part] (C) -- ++(-120:1.169) node[below left] {\small $mg\cos\ang{30}$};
   % normal force and friction
   \draw[force, accent] (C) -- ++(60:1.169) node[above right]
      {\small $F_{\mathrm{n}}$};
   \draw[force, oxblood] ($(C)+(150:0.45)$) -- ++(150:0.234)
      node[above left] {\small $F_{\mathrm{fr}}$};
   \draw[black!50, line width=0.4pt] ($(C)+(-90:0.5)$) arc[start angle=-90,
      end angle=-120, radius=0.5];
   \node at ($(C)+(-105:0.72)$) {\scriptsize \ang{30}};
   \fill (C) circle (1.3pt);
\end{tikzpicture}
\caption{A block sliding down a rough ramp. Its weight $mg$ is split
(dashed) into a part along the ramp and a part into it. The normal force
$F_{\mathrm{n}}$ (teal) balances the part into the ramp, and friction
$F_{\mathrm{fr}}$ (dark red) acts up the ramp, against the motion.}
\label{fig:en-ramp}
\end{figure}

Over the length $L = \SI{3}{\metre}$ the block drops by $L\sin\ang{30} =
\SI{1.5}{\metre}$, and the three forces do the work
\begin{align*}
   W_{\mathrm{weight}} &= -mg\times(-1.5)
      && \text{by \cref{sec:en-track}, $-mg$ times the rise,}\\
   &= 2\times9.80665\times1.5 = \SI{29.42}{\joule} ,\\
   W_{\mathrm{n}} &= 0
      && \text{perpendicular to the motion,}\\
   W_{\mathrm{fr}} &= -3.397\times3 = \SI{-10.19}{\joule}
      && \text{against the motion.}
\end{align*}
The block arrives with $K = 29.42 - 10.19 = \SI{19.23}{\joule}$, and so
with speed $\sqrt{2\times19.23/2} = \SI{4.385}{\metre\per\second}$; on a
smooth ramp it would arrive at $\sqrt{2g\times1.5} =
\SI{5.424}{\metre\per\second}$.

\subsection*{The rate of energy loss}

Friction and the drag of \cref{sec:ne-drag} always point against the
velocity, so on a trip out to a point and back they do negative work on
both legs, whereas a force with a potential energy does the work $U(x_A)
- U(x_A) = 0$ on any trip that returns to its start $x_A$. They therefore
have no potential energy. To see their effect on $E$, write the net force
on a body moving along a line as a conservative part, $-\dd U/\dd x$,
plus the rest, $F^{\mathrm{nc}}$, the non-conservative part. The kinetic
energy changes at the rate $(-\dd U/\dd x + F^{\mathrm{nc}})\,v$, by
\cref{sec:en-varying}, and the potential energy at the rate $(\dd U/\dd
x)\,v$, by the chain rule, so their sum changes at the rate
\begin{equation}
   \frac{\dd E}{\dd t} = F^{\mathrm{nc}}\,v :
   \label{eq:en-dEdt-1d}
\end{equation}
the total energy changes at exactly the rate at which the
non-conservative force does work. Friction always opposes the velocity,
so $F^{\mathrm{nc}}v = -\mu_{\mathrm{k}}F_{\mathrm{n}}\lvert v\rvert$, and
the energy can only fall. For linear drag, $F^{\mathrm{nc}} = -bv$ and
$\dd E/\dd t = -bv^2$.

The bead of \cref{sec:ne-drag} sinking at its terminal velocity shows
this rate at work. Measuring the depth $x$ downwards, its potential
energy is $U = -mgx$, which falls at the rate $mg\,v_\infty$, while its
kinetic energy stays constant. Every second, then, the weight does
$mg\,v_\infty = 0.01\times9.80665\times1.961 = \SI{0.1923}{\joule}$ of
work, and the drag removes all of it: $bv_\infty^2 = 0.05\times1.961^2 =
\SI{0.1923}{\watt}$, the same, because $bv_\infty = mg$.

\begin{selfcheck}
A box sliding across a floor stops after \SI{0.68}{\metre}. How far would
a box of twice the mass, sliding on the same floor at the same starting
speed, travel before stopping? Where has the kinetic energy gone?
\end{selfcheck}
